\section{Methods}
\label{sec:methods}

\subsection{Protocol}
Every model predicts $\log_{10}$(cycle life) from features of cycles $\le N$
(Table~\ref{tab:featsets}). Missing values are filled with the training
median and features are standardised on the training cells. To keep the
small test sets honest, all choices (candidate features, model variants,
settings and the bounds of the GP) are made on the 41 training cells only,
by cross-validation or, where a choice is itself being compared, nested
cross-validation. Each configuration chosen this way is then evaluated once
on the primary and secondary test sets, and all are reported. The unseen
batch u1 is used once, at the end, with every model frozen
(Section~\ref{sec:rq5}). Models using the
temperature integral leave out the four cells with a detached probe, as in
the paper. Randomised steps are repeated with 5--10 seeds; the spread is a few
cycles (at most 11) and is not shown.

\subsection{Metrics}
Predictions of $\log_{10}$ life are converted back to cycles before
scoring. Table~\ref{tab:metrics} lists the scores; RMSE and MAPE are those
of~\cite{severson2019}, the last three are ours.

\begin{table}[!htb]
\centering
\caption{Evaluation metrics. $y_i$, $\hat y_i$: observed and predicted
life of test cell $i$ (cycles).}
\label{tab:metrics}
\small
\setlength{\tabcolsep}{2.5pt}
\begin{tabular}{@{}>{\raggedright\arraybackslash}p{0.20\columnwidth}>{\raggedright\arraybackslash}p{0.42\columnwidth}>{\raggedright\arraybackslash}p{0.34\columnwidth}@{}}
\toprule
Metric & Definition & Purpose / values \\
\midrule
RMSE (cycles) & $\sqrt{\tfrac1n\sum_i(\hat y_i-y_i)^2}$ & typical error; dominated by large errors \\
MAPE (\%) & $\tfrac{100}{n}\sum_i|\hat y_i-y_i|/y_i$ & error relative to life; fair to short-lived cells \\
Constant baseline & every cell predicted at the geometric mean of training lives, 622 cycles (mean of $\log_{10}$ life) & the bar to beat: RMSE 401 (primary), 511 (secondary) \\
In-range RMSE & RMSE over test cells with life in the training range, 300--2160 & separates prediction from extrapolation; drops \texttt{b1c2} (2237) and \texttt{b2c1} (148) \\
Paired bootstrap & RMSE difference of two models on the same resampled test cells; 5000 resamples, 95\,\% interval & difference is real only if the interval excludes 0 \\
GP coverage & share of test cells inside the GP's 95\,\% prediction interval & calibrated if $\approx95$\,\% \\
\bottomrule
\end{tabular}
\end{table}

\subsection{Models}
Table~\ref{tab:models} lists the models; Section~\ref{sec:rq2} explains
the non-linear ones. Trees and neural networks were not among the main
models: a tree cannot predict beyond the training lives, and a neural
network has far more parameters than the 41 cells. One random forest was
tested as a check.

\begin{table}[!htb]
\centering
\caption{Models. All use standardised inputs and predict $\log_{10}$ life;
CV: 4-fold cross-validation repeated 5--10 times on the training cells.}
\label{tab:models}
\small
\setlength{\tabcolsep}{2.5pt}
\begin{tabular}{@{}>{\raggedright\arraybackslash}p{0.14\columnwidth}>{\raggedright\arraybackslash}p{0.46\columnwidth}>{\raggedright\arraybackslash}p{0.35\columnwidth}@{}}
\toprule
Model & Setup & Tuning \\
\midrule
EN~\cite{severson2019} & linear; $\ell_1+\ell_2$ penalty, $\alpha\|w\|_1\rho+\tfrac{\alpha}{2}(1-\rho)\|w\|_2^2$; plain least squares for F1 & $\alpha$ (100 values), $\rho\in[0.1,1]$ (8); CV \\
GP & kernel: linear $+$ Mat\'ern-5/2 (one length-scale per feature) $+$ noise; gives a predictive s.d. & marginal likelihood; length-scale $\ge1$, noise $\ge0.01$ \\
GAM & cubic spline per feature (linear beyond the data) $+$ ridge & knots (3--5), ridge (11); CV \\
RF & 300 trees on bootstrap samples (check only) & leaf size, features per split; CV \\
\bottomrule
\end{tabular}
\end{table}

\subsection{Experiments}
Table~\ref{tab:experiments} lists which models and feature sets answer each
research question. The non-linear models are run on F2 only: it is the
paper's best set on a new batch, while F3 contains batch-dependent features
(Section~\ref{sec:eda}) that already make the elastic net fail on the
secondary batch. No new feature set is used, because none of the screened
candidates added information (Section~\ref{sec:rq3}).

\begin{table}[!htb]
\centering
\caption{Models and feature sets used for each research question. Prim./Sec.:
primary and secondary test sets.}
\label{tab:experiments}
\small
\setlength{\tabcolsep}{2.5pt}
\begin{tabular}{@{}l>{\raggedright\arraybackslash}p{0.58\columnwidth}>{\raggedright\arraybackslash}p{0.27\columnwidth}@{}}
\toprule
RQ & Models $\times$ feature sets & Tested on \\
\midrule
RQ1 & EN on F1 (plain least squares), F2, F3 & Prim., Sec. \\
RQ3 & 19 candidates screened with linear fits; overhang check: EN on F1 and F2 with and without it & Training; Prim., Sec., u1 \\
RQ2 & GP, GAM and a random forest on F2, against EN-F2 & Prim., Sec. \\
RQ4 & EN on F1, F2, F3, recomputed for $N=20$--150 & Prim., Sec. \\
RQ5 & EN on F1, F2, F3; GP, GAM and random forest on F2 & u1 standard and \texttt{newstructure} \\
\bottomrule
\end{tabular}
\end{table}
